\begin{proposition}
\label{proposition-homotopy-category-triangulated}
Let $\mathcal{A}$ be an additive category.
The category $K(\mathcal{A})$ of complexes up to
homotopy with its natural translation functors
and distinguished triangles as defined above
is a triangulated category.
\end{proposition}

\begin{proof}
Proof of TR1. By definition every triangle isomorphic to a distinguished
one is distinguished. Also, any triangle $(A^\bullet, A^\bullet, 0, 1, 0, 0)$
is distinguished since $0 \to A^\bullet \to A^\bullet \to 0 \to 0$ is
a termwise split sequence of complexes. Finally, given any morphism of
complexes $f : K^\bullet \to L^\bullet$ the triangle
$(K, L, C(f), f, i, -p)$ is distinguished by
Lemma \ref{lemma-the-same-up-to-isomorphisms}.

\medskip\noindent
Proof of TR2. Let $(X, Y, Z, f, g, h)$ be a triangle.
Assume $(Y, Z, X[1], g, h, -f[1])$ is distinguished.
Then there exists a termwise split sequence of complexes
$A^\bullet \to B^\bullet \to C^\bullet$ such that the associated
triangle $(A^\bullet, B^\bullet, C^\bullet, \alpha, \beta, \delta)$
is isomorphic to $(Y, Z, X[1], g, h, -f[1])$. Rotating back we see
that $(X, Y, Z, f, g, h)$ is isomorphic to
$(C^\bullet[-1], A^\bullet, B^\bullet, -\delta[-1], \alpha, \beta)$.
It follows from Lemma \ref{lemma-rotate-triangle} that the triangle
$(C^\bullet[-1], A^\bullet, B^\bullet, \delta[-1], \alpha, \beta)$
is isomorphic to
$(C^\bullet[-1], A^\bullet, C(\delta[-1])^\bullet, \delta[-1], i, p)$.
Precomposing the previous isomorphism of triangles with $-1$ on $Y$
it follows that $(X, Y, Z, f, g, h)$ is isomorphic to
$(C^\bullet[-1], A^\bullet, C(\delta[-1])^\bullet, \delta[-1], i, -p)$.
Hence it is distinguished by
Lemma \ref{lemma-the-same-up-to-isomorphisms}.
On the other hand, suppose that $(X, Y, Z, f, g, h)$ is distinguished.
By Lemma \ref{lemma-the-same-up-to-isomorphisms} this means that it is
isomorphic to a triangle of the form
$(K^\bullet, L^\bullet, C(f), f, i, -p)$ for some morphism of
complexes $f$. Then the rotated triangle $(Y, Z, X[1], g, h, -f[1])$ is
isomorphic to $(L^\bullet, C(f), K^\bullet[1], i, -p, -f[1])$ which is
isomorphic to the triangle $(L^\bullet, C(f), K^\bullet[1], i, p, f[1])$.
By Lemma \ref{lemma-rotate-cone} this triangle is distinguished.
Hence $(Y, Z, X[1], g, h, -f[1])$ is distinguished as desired.

\medskip\noindent
Proof of TR3. Let
$(X, Y, Z, f, g, h)$ and $(X', Y', Z', f', g', h')$
be distinguished triangles of $K(\mathcal{A})$
and let $a : X \to X'$ and $b : Y \to Y'$ be morphisms
such that $f' \circ a = b \circ f$. By
Lemma \ref{lemma-the-same-up-to-isomorphisms} we may assume that
$(X, Y, Z, f, g, h) = (X, Y, C(f), f, i, -p)$ and
$(X', Y', Z', f', g', h') = (X', Y', C(f'), f', i', -p')$.
At this point we simply apply Lemma \ref{lemma-functorial-cone}
to the commutative diagram given by $f, f', a, b$.

\medskip\noindent
Proof of TR4. At this point we know that $K(\mathcal{A})$
is a pre-triangulated category. Hence we can use
Lemma \ref{lemma-easier-axiom-four}. Let $A^\bullet \to B^\bullet$
and $B^\bullet \to C^\bullet$ be composable morphisms of
$K(\mathcal{A})$. By Lemma \ref{lemma-sequence-maps-split} we may assume that
$A^\bullet \to B^\bullet$ and $B^\bullet \to C^\bullet$
are split injective morphisms. In this case the result follows
from Lemma \ref{lemma-two-split-injections}.
\end{proof}
